% =============================================================================
% Informe de Solución — Ejercicios de Análisis de Componentes Principales
% Fuente: Ayudantía Sebastián Chávez Orellana, IC USACH, 1.er sem. 2026
% Prof. Max Chacón Pacheco
% =============================================================================
\documentclass[11pt,a4paper]{article}

\usepackage[utf8]{inputenc}
\usepackage[T1]{fontenc}
\usepackage[spanish,es-tabla]{babel}
\usepackage[margin=2.2cm]{geometry}
\usepackage{amsmath,amssymb,amsthm,mathtools,bm}
\usepackage{booktabs,array}
\usepackage{xcolor}
\usepackage{tcolorbox}
\tcbuselibrary{skins,breakable}
\usepackage{hyperref}
\usepackage{enumitem}
\usepackage{tikz}
\usetikzlibrary{arrows.meta,positioning,calc,patterns,decorations.markings,shapes.geometric}
\usepackage{pgfplots}
\pgfplotsset{compat=1.18}
\usepackage{fancyhdr}
\usepackage{titlesec}
\usepackage{caption}
\usepackage{float}

\definecolor{usachteal}{HTML}{009999}
\definecolor{usachdark}{HTML}{003333}
\hypersetup{colorlinks=true,linkcolor=teal!70!black,urlcolor=blue!70!black}

\newtcolorbox{enuncbox}[1][Enunciado]{enhanced,breakable,
  colback=blue!3,colframe=blue!50!black,
  fonttitle=\bfseries\color{white},coltitle=white,
  colbacktitle=blue!60!black,title={#1},
  sharp corners=south,boxrule=0.6pt}

\newtcolorbox{solbox}[1][Resolución]{enhanced,breakable,
  colback=usachteal!5,colframe=usachteal!70!black,
  fonttitle=\bfseries\color{white},coltitle=white,
  colbacktitle=usachteal!80!black,title={#1},
  sharp corners=south,boxrule=0.6pt}

\newtcolorbox{respbox}[1][Respuesta]{enhanced,breakable,
  colback=orange!8,colframe=orange!70!black,
  fonttitle=\bfseries,title={#1},
  sharp corners=south,boxrule=0.5pt}

\pagestyle{fancy}
\fancyhf{}
\fancyhead[L]{\small\textsc{Inteligencia Computacional · USACH}}
\fancyhead[R]{\small\textsc{Informe ACP — Ayudantía 2026/I}}
\fancyfoot[C]{\thepage}
\renewcommand{\headrulewidth}{0.4pt}

\titleformat{\section}{\Large\bfseries\color{usachdark}}{\thesection.}{0.6em}{}
\titleformat{\subsection}{\large\bfseries\color{usachteal!80!black}}{\thesubsection.}{0.5em}{}

% Macros para biplots ACP (flechas de variables y puntos etiquetados)
\newcommand{\arr}[4]{%
  \draw[-{Stealth[length=2mm]},thick,usachteal!80!black] (axis cs:0,0) -- (axis cs:#1,#2);%
  \node[anchor=#4,font=\small] at (axis cs:#1,#2) {\textbf{#3}};%
}
\newcommand{\pt}[3]{%
  \addplot[only marks,mark=diamond*,mark size=3pt,color=blue] coordinates {(#1,#2)};%
  \node[anchor=south west,font=\small\bfseries] at (axis cs:#1,#2) {#3};%
}
\newcommand{\ptt}[3]{%
  \addplot[only marks,mark=diamond*,mark size=3pt,color=usachteal!80!black] coordinates {(#1,#2)};%
  \node[anchor=south west,font=\small\bfseries] at (axis cs:#1,#2) {#3};%
}
\newcommand{\pto}[3]{%
  \addplot[only marks,mark=diamond*,mark size=3pt,color=orange!80!black] coordinates {(#1,#2)};%
  \node[anchor=south west,font=\small\bfseries] at (axis cs:#1,#2) {#3};%
}
\newcommand{\ptn}[3]{%
  \addplot[only marks,mark=*,mark size=3pt,color=red!70!black] coordinates {(#1,#2)};%
  \node[anchor=south west,font=\small\bfseries] at (axis cs:#1,#2) {#3};%
}

\begin{document}

\begin{titlepage}
\centering
\vspace*{2cm}
{\Huge\bfseries\color{usachdark} Informe de Solución\\[0.3cm]}
{\LARGE\color{usachteal} Ejercicios de Análisis de\\Componentes Principales}\\[1cm]
{\large \textit{Guía ayudantía — Primer semestre 2026}}\\[2cm]
\begin{tcolorbox}[width=0.75\textwidth,colback=usachteal!10,colframe=usachteal!70!black]
\centering
\textbf{Inteligencia Computacional}\\
Magíster en Ingeniería\\
\textbf{Universidad de Santiago de Chile}\\[0.3cm]
Prof.\ Max Chacón Pacheco\\
Ayudante: Sebastián Chávez Orellana
\end{tcolorbox}
\vspace{1.5cm}

{\small Resolución paso a paso de los 7 problemas con gráficos generados en TikZ/PGFPlots.}\\[0.3cm]
{\small\today}
\vfill
\end{titlepage}

\tableofcontents
\newpage

% =============================================================================
\section{Marco general y convenciones}
% =============================================================================

Sea $X\in\mathbb{R}^{n\times p}$ la matriz de datos estandarizada (media cero, desviación estándar uno).
La matriz de correlación $R = \tfrac{1}{n-1} X^\top X$ es simétrica y semidefinida positiva, y se descompone en
\[
R \;=\; V\,\Lambda\,V^\top, \qquad V = [\,\bm v_1\,|\,\bm v_2\,|\cdots|\,\bm v_p\,],\quad \Lambda = \mathrm{diag}(\lambda_1,\ldots,\lambda_p).
\]
Cada $\lambda_i$ es la \textbf{varianza explicada} por la componente $i$, y $\bm v_i$ es su \textbf{dirección} (autovector). Las \emph{cargas factoriales} son $\ell_{ji}=v_{ji}\sqrt{\lambda_i}$ y coinciden con la correlación $\mathrm{corr}(X_j,\mathrm{CP}_i)$.

\begin{tcolorbox}[colback=yellow!10,colframe=yellow!50!black]
\textbf{Fórmulas clave.}
\begin{align*}
\sum_{i=1}^p \lambda_i &= p \qquad\text{(traza de }R\text{, datos estandarizados)}\\[2pt]
\text{Varianza explicada con $k$ componentes:}\quad \mathrm{VE}_k &= \frac{\sum_{i=1}^k \lambda_i}{\sum_{i=1}^p \lambda_i}\cdot 100\%\\[2pt]
\text{Información perdida:}\quad \mathrm{IP}_k &= 100\% - \mathrm{VE}_k\\[2pt]
\text{Regla de Kaiser:}\quad &\text{retener CP con }\lambda_i>1.
\end{align*}
\end{tcolorbox}

\noindent En todos los problemas se usa la convención: \emph{columnas} de la tabla de autovectores = direcciones $\bm v_i$ (coordenadas $v_{ji}$ por variable $j$).

\newpage
% =============================================================================
\section{P1. Comportamiento de consumo alimenticio}
% =============================================================================

\begin{enuncbox}
Familias francesas clasificadas por su jefe de hogar: \textbf{MA\#} (manual), \textbf{EM\#} (empleado), \textbf{PF\#} (profesional), con \# $\in\{2,3,4,5\}$ número de hijos. Se mide el consumo promedio de 7 productos.

\medskip
Autovalores: $\lambda=(4{,}339;\;1{,}829;\;0{,}625;\;0{,}502;\;0{,}393;\;0{,}099;\;0{,}083)$.

\smallskip
Autovectores de las dos primeras componentes:

\begin{center}\small
\begin{tabular}{lrr}\toprule
Variable & $\bm v_1$ & $\bm v_2$ \\\midrule
Pan      & $-0{,}497$ & $\phantom{-}0{,}841$ \\
Verduras & $-0{,}972$ & $\phantom{-}0{,}131$ \\
Frutas   & $-0{,}931$ & $-0{,}277$ \\
Carnes   & $-0{,}963$ & $-0{,}190$ \\
Aves     & $-0{,}912$ & $-0{,}265$ \\
Lácteos  & $-0{,}584$ & $\phantom{-}0{,}707$ \\
Vinos    & $\phantom{-}0{,}425$ & $\phantom{-}0{,}649$ \\\bottomrule
\end{tabular}
\end{center}
\end{enuncbox}

\begin{solbox}
\textbf{(a) Información perdida con CP1 + CP2.}

Suma de autovalores ($\approx p = 7$, los decimales son por redondeo en la guía):
\[
\sum_{i=1}^{7}\lambda_i \;=\; 4{,}339+1{,}829+0{,}625+0{,}502+0{,}393+0{,}099+0{,}083 \;=\; 7{,}870.
\]
Varianza explicada por las dos primeras:
\[
\mathrm{VE}_2 \;=\; \frac{4{,}339+1{,}829}{7{,}870}\cdot 100\% \;=\; \frac{6{,}168}{7{,}870}\cdot 100\% \;\approx\; 78{,}37\%.
\]
\begin{respbox}
$\boxed{\mathrm{IP}_2 \approx 21{,}63\%}$ de información se pierde al proyectar a dos componentes.
\end{respbox}

\textbf{(b) Caracterización inspeccionando la nube de familias.}

Del diagrama original (fig.~\ref{fig:p1nube}) se observa que los \emph{profesionales} (PF) quedan a la izquierda y abajo, los \emph{empleados} (EM) y \emph{manuales} (MA) se distribuyen mayormente a la derecha y arriba. Como CP1 ordena por ingreso (PF a la izquierda con CP1$<0$) y CP2 separa por tamaño/composición familiar:

\begin{itemize}[noitemsep,leftmargin=*]
\item \textbf{CP1} $\equiv$ \emph{nivel de ingreso / poder adquisitivo}: CP1$<0$ $\Rightarrow$ mayor ingreso (PF); CP1$>0$ $\Rightarrow$ menor ingreso (MA).
\item \textbf{CP2} $\equiv$ \emph{tamaño familiar / dieta básica}: CP2$>0$ $\Rightarrow$ más hijos, mayor consumo de pan y lácteos.
\end{itemize}

\textbf{(c) Variables originales en el plano.} Se grafican los puntos $(v_{j1},v_{j2})$ por variable (fig.~\ref{fig:p1var}).

\textbf{(d) Asociación componente $\leftrightarrow$ productos.}
\begin{itemize}[noitemsep,leftmargin=*]
\item \textbf{CP1}: dominada por Verduras, Frutas, Carnes, Aves, Lácteos (todas con $v_{j1}\approx-0{,}9$) frente a Vinos (único positivo). Es el \textbf{eje de la canasta de "calidad" alimenticia}: CP1 muy negativo $\Rightarrow$ alto consumo proteico-vegetal; CP1 positivo $\Rightarrow$ predominio del vino.
\item \textbf{CP2}: contrapone \textbf{Pan + Lácteos} (cargas $+0{,}84$ y $+0{,}71$) frente a \textbf{Carnes, Aves, Frutas} ($<0$). Es el \textbf{eje del estilo dietético}: CP2$>0$ $\Rightarrow$ dieta de hidratos+lácteos (familias con niños); CP2$<0$ $\Rightarrow$ dieta proteico-vegetal.
\end{itemize}

\textbf{(e) Relación productos $\leftrightarrow$ condición socioeconómica.}

\begin{itemize}[noitemsep,leftmargin=*]
\item \textbf{PF} (cuadrante III, CP1$<0$, CP2$<0$): alto consumo de carnes, aves, frutas, verduras (canasta cara).
\item \textbf{EM} (cuadrante II, CP1$<0$, CP2$>0$): mayor peso de pan y lácteos $\Rightarrow$ familias con niños.
\item \textbf{MA} (cuadrantes I y IV, CP1$>0$): consumo más asociado al vino y dieta restringida.
\end{itemize}
\end{solbox}

\begin{figure}[H]
\centering
\begin{tikzpicture}[scale=1]
\begin{axis}[
  width=11.5cm, height=8cm,
  axis lines=middle, xlabel={CP1}, ylabel={CP2},
  xlabel near ticks, ylabel near ticks,
  xmin=-3,xmax=3,ymin=-2.5,ymax=2.5,
  grid=both, grid style={gray!20},
  title={Familias en el plano CP1--CP2 (reconstrucción cualitativa)}]
% PF (profesionales) - cuadrante III
\addplot[only marks,mark=*,mark size=2pt,color=blue]
  coordinates {(-2.3,-0.1) (-1.1,-1.0) (-0.8,-1.2)};
\node[blue,anchor=south east] at (axis cs:-2.3,-0.1) {\small PF};
\node[blue,anchor=west] at (axis cs:-1.0,-1.0) {\small PF};
\node[blue,anchor=west] at (axis cs:-0.7,-1.3) {\small P};
% EM (empleados) - mayormente cuadrante II
\addplot[only marks,mark=square*,mark size=2pt,color=red]
  coordinates {(-1.9,1.2) (0.6,0.9) (1.5,-0.1) (1.8,-1.6)};
\node[red,anchor=south west] at (axis cs:-1.9,1.2) {\small EM};
\node[red,anchor=south] at (axis cs:0.6,0.9) {\small EM};
\node[red,anchor=south] at (axis cs:1.5,-0.1) {\small EM};
\node[red,anchor=west] at (axis cs:1.8,-1.6) {\small EM};
% MA (manuales) - I y eje
\addplot[only marks,mark=triangle*,mark size=2.5pt,color=usachteal]
  coordinates {(-0.05,2.2) (1.1,1.1) (2.0,0.5) (2.3,-0.05) (0.7,-0.6)};
\node[usachteal,anchor=south east] at (axis cs:-0.05,2.2) {\small MA};
\node[usachteal,anchor=south] at (axis cs:1.1,1.1) {\small MA};
\node[usachteal,anchor=south] at (axis cs:2.0,0.5) {\small MA};
\node[usachteal,anchor=west] at (axis cs:2.3,-0.05) {\small MA};
\node[usachteal,anchor=west] at (axis cs:0.7,-0.6) {\small PF};
\end{axis}
\end{tikzpicture}
\caption{Distribución de familias en el plano CP1--CP2 (reconstrucción a partir de la figura del enunciado).}
\label{fig:p1nube}
\end{figure}

\begin{figure}[H]
\centering
\begin{tikzpicture}
\begin{axis}[
  width=11.5cm, height=8cm,
  axis lines=middle, xlabel={$v_{j1}$ (CP1)}, ylabel={$v_{j2}$ (CP2)},
  xlabel near ticks, ylabel near ticks,
  xmin=-1.1,xmax=1.1,ymin=-1.1,ymax=1.1,
  grid=both, grid style={gray!20},
  title={Variables originales (autovectores) y círculo unitario}]
% Círculo de correlaciones
\draw[gray, dashed] (axis cs:0,0) circle [radius=1];
% Vectores (desenrollados)
\def\arr#1#2#3#4{\draw[-{Stealth[length=2mm]},thick,usachteal!80!black] (axis cs:0,0) -- (axis cs:#1,#2); \node[anchor=#4,font=\small] at (axis cs:#1,#2) {\textbf{#3}};}
\arr{-0.497}{0.841}{Pan}{south west}
\arr{-0.972}{0.131}{Verduras}{south}
\arr{-0.931}{-0.277}{Frutas}{north}
\arr{-0.963}{-0.190}{Carnes}{south}
\arr{-0.912}{-0.265}{Aves}{west}
\arr{-0.584}{0.707}{Lácteos}{south east}
\arr{0.425}{0.649}{Vinos}{south west}
\end{axis}
\end{tikzpicture}
\caption{Variables originales del problema P1 proyectadas sobre el plano CP1--CP2.}
\label{fig:p1var}
\end{figure}

\newpage
% =============================================================================
\section{P2. Caracterización de billetes falsos}
% =============================================================================

\begin{enuncbox}
6 variables sobre billetes suizos: \textbf{LON} (longitud), \textbf{LD} (diagonal), \textbf{AI} (ancho izq), \textbf{AD} (ancho der), \textbf{AMI} (margen inf), \textbf{AMS} (margen sup). Tres grupos: papel originales ($\bullet$), plástico originales ($\blacksquare$), falsos ($\blacktriangle$).

\smallskip
$\lambda=(2{,}58;\;1{,}34;\;0{,}76;\;0{,}56;\;0{,}50;\;0{,}26)$ \quad (suman $6{,}00=p$).

\begin{center}\small
\begin{tabular}{lrr}\toprule
Variable & $\bm v_1$ & $\bm v_2$ \\\midrule
LON & $0{,}395$ & $\phantom{-}0{,}799$\\
LD  & $0{,}207$ & $\phantom{-}0{,}345$\\
AI  & $0{,}445$ & $-0{,}263$\\
AD  & $0{,}411$ & $-0{,}375$\\
AMI & $0{,}347$ & $-0{,}072$\\
AMS & $0{,}560$ & $-0{,}163$\\\bottomrule
\end{tabular}
\end{center}

\smallskip
Adicionalmente, sobre monedas (6 medidas): $\lambda=(1{,}96;\;1{,}54;\;1{,}09;\;0{,}73;\;0{,}40;\;0{,}28)$.
\end{enuncbox}

\begin{solbox}
\textbf{(a) Validez del análisis (billetes).}
\[
\mathrm{VE}_2 = \frac{2{,}58+1{,}34}{6{,}00}\cdot 100\% = \frac{3{,}92}{6}\cdot 100\% \approx \boxed{65{,}33\%}.
\]

\textbf{(b) Interpretación de las componentes.}
\begin{itemize}[noitemsep,leftmargin=*]
\item \textbf{CP1}: todas las cargas son \emph{positivas} y de magnitud similar ($0{,}21$--$0{,}56$). Las más altas son \textbf{AMS} (0,56), \textbf{AI} (0,45) y \textbf{AD} (0,41). $\Rightarrow$ \textbf{CP1 mide la "forma" del billete: dimensión transversal y márgenes internos.}
\item \textbf{CP2}: signo positivo en \textbf{LON} (0,80) y \textbf{LD} (0,35); negativo en los anchos y márgenes. Contrapone tamaño longitudinal vs.\ ancho. $\Rightarrow$ \textbf{CP2 mide el tamaño longitudinal/formato.}
\end{itemize}

Para graficar las variables en el círculo de correlaciones no se usan directamente los autovectores $v_{ji}$, sino las \textbf{cargas factoriales}
\[
\ell_{ji}=v_{ji}\sqrt{\lambda_i}.
\]
En P2, $\sqrt{\lambda_1}=\sqrt{2{,}58}\approx1{,}606$ y $\sqrt{\lambda_2}=\sqrt{1{,}34}\approx1{,}158$. Por tanto, por ejemplo, LON queda en
\[
(\ell_{\mathrm{LON},1},\ell_{\mathrm{LON},2})=(0{,}395\sqrt{2{,}58},\;0{,}799\sqrt{1{,}34})\approx(0{,}634,\;0{,}925).
\]

\textbf{(c) Características de los billetes originales.} De la figura del enunciado:
\begin{itemize}[noitemsep,leftmargin=*]
\item \textbf{Papel originales} ($\bullet$, cuadrante I): CP1$>0$ y CP2$>0$ $\Rightarrow$ márgenes amplios \emph{y} billete más largo.
\item \textbf{Plástico originales} ($\blacksquare$, cuadrante IV): CP1$>0$ y CP2$<0$ $\Rightarrow$ comparten márgenes amplios con los de papel pero son más cortos.
\end{itemize}
La característica común de \emph{todos} los originales es \textbf{CP1 alto}: \textbf{márgenes amplios}.

\textbf{(d) Diferencias entre falsificaciones.} Los falsos se ubican en CP1$<0$ (márgenes estrechos) y se separan en \textbf{dos sub-grupos}:
\begin{itemize}[noitemsep,leftmargin=*]
\item Cuadrante II (CP2$>0$): falsos que imitan a los originales de papel.
\item Cuadrante III (CP2$<0$): falsos que imitan a los originales de plástico.
\end{itemize}
$\Rightarrow$ Sí existen diferencias claras: hay \textbf{dos tipos de falsificaciones}.

\textbf{(e) Comparación con monedas.}
\[
\mathrm{VE}_2^{\text{monedas}} = \frac{1{,}96+1{,}54}{6}\cdot 100\% = \frac{3{,}50}{6}\cdot 100\% \approx \boxed{58{,}33\%}.
\]
Se pierde $\approx 7$ p.p.\ respecto a los billetes. La regla de Kaiser indica que en monedas las \emph{tres} primeras componentes tienen $\lambda>1$; conviene retener tres ($\mathrm{VE}_3 = 4{,}59/6 \approx 76{,}5\%$) para alcanzar precisión equivalente.
\end{solbox}

\begin{figure}[H]
\centering
\begin{tikzpicture}
\begin{axis}[
  width=11cm, height=8cm,
  axis lines=middle, xlabel={CP1}, ylabel={CP2},
  xmin=-3,xmax=3,ymin=-2.5,ymax=2.5,
  grid=both, grid style={gray!20},
  title={Distribución cualitativa de billetes (P2)}]
% Papel (Q1)
\addplot[only marks,mark=*,mark size=2pt,color=blue]
  coordinates {(1.2,1.6) (1.4,1.4) (1.6,1.7) (1.0,1.2) (1.8,1.5) (1.3,1.9) (1.5,1.3) (1.7,1.1) (1.1,1.5)};
\node[blue,anchor=south] at (axis cs:1.4,2.0) {\textbf{Papel}};
% Plástico (Q4)
\addplot[only marks,mark=square*,mark size=2pt,color=red]
  coordinates {(1.2,-1.0) (1.5,-1.4) (1.8,-1.1) (1.4,-0.9) (1.6,-1.3) (1.9,-1.0) (1.3,-1.5)};
\node[red,anchor=north] at (axis cs:1.5,-1.7) {\textbf{Plástico}};
% Falsos sub-grupo II
\addplot[only marks,mark=triangle*,mark size=2.5pt,color=usachteal!80!black]
  coordinates {(-1.5,1.0) (-1.8,1.2) (-2.0,0.8) (-1.6,0.6) (-2.2,1.1) (-1.7,0.9)};
% Falsos sub-grupo III
\addplot[only marks,mark=triangle*,mark size=2.5pt,color=usachteal!80!black]
  coordinates {(-1.6,-0.6) (-1.9,-0.4) (-2.1,-0.8) (-1.7,-1.0) (-2.3,-0.5) (-1.8,-1.2)};
\node[usachteal!80!black,anchor=east] at (axis cs:-2.4,0.3) {\textbf{Falsos}};
\end{axis}
\end{tikzpicture}
\caption{Reconstrucción de la dispersión de billetes en CP1--CP2 (P2).}
\end{figure}

\begin{figure}[H]
\centering
\begin{tikzpicture}
\begin{axis}[
  width=11cm,height=8cm, axis lines=middle, xlabel={$\ell_{j1}$ (CP1)}, ylabel={$\ell_{j2}$ (CP2)},
  xmin=-1.1,xmax=1.1,ymin=-1.1,ymax=1.1, grid=both, grid style={gray!20},
  title={Círculo de correlaciones P2 (billetes)}]
\draw[gray, dashed] (axis cs:0,0) circle [radius=1];
\arr{0.634}{0.925}{LON}{south west}
\arr{0.332}{0.399}{LD}{south west}
\arr{0.715}{-0.304}{AI}{north}
\arr{0.660}{-0.434}{AD}{north west}
\arr{0.557}{-0.083}{AMI}{north}
\arr{0.899}{-0.189}{AMS}{east}
\end{axis}
\end{tikzpicture}
\caption{Cargas factoriales del problema P2: $\ell_{ji}=v_{ji}\sqrt{\lambda_i}$, no autovectores crudos.}
\end{figure}

\newpage
% =============================================================================
\section{P3. Caracterización de servicios hospitalarios}
% =============================================================================

\begin{enuncbox}
Hospital de Andalucía, 22.846 ingresos, 13 servicios, 7 variables: \textbf{NI} (n.\ ingresos), \textbf{MO} (mortalidad), \textbf{RE} (reingresos), \textbf{NE} (consultas externas), \textbf{ICM} (complejidad), \textbf{ES} (estancias), \textbf{IF} (funcionalidad).

\smallskip
$\lambda_1 = 2{,}558$,\quad $\lambda_2 = 1{,}829$ (datos estandarizados, $\sum\lambda_i=p=7$).

\begin{center}\small
\begin{tabular}{lrr}\toprule
Variable & $\bm v_1$ & $\bm v_2$\\\midrule
NI  & $\phantom{-}0{,}860$ & $-0{,}066$\\
MO  & $\phantom{-}0{,}421$ & $\phantom{-}0{,}747$\\
RE  & $-0{,}406$ & $\phantom{-}0{,}670$\\
NE  & $-0{,}250$ & $\phantom{-}0{,}388$\\
ICM & $-0{,}562$ & $\phantom{-}0{,}635$\\
ES  & $\phantom{-}0{,}820$ & $\phantom{-}0{,}508$\\
IF  & $\phantom{-}0{,}663$ & $\phantom{-}0{,}078$\\\bottomrule
\end{tabular}
\end{center}
\end{enuncbox}

\begin{solbox}
\textbf{(a) Validez.}
\[
\mathrm{VE}_2 \;=\; \frac{2{,}558+1{,}829}{7}\cdot 100\% \;=\; \frac{4{,}387}{7}\cdot 100\% \;\approx\; \boxed{62{,}67\%}.
\]

\textbf{(b)} Variables ubicadas en $(v_{j1},v_{j2})$ — ver figura~\ref{fig:p3var}.

\textbf{(c) Interpretación.}
\begin{itemize}[noitemsep,leftmargin=*]
\item \textbf{CP1 (Demanda / volumen del servicio):} cargas positivas en NI ($0{,}86$), ES ($0{,}82$), IF ($0{,}66$); negativas en ICM, RE, NE. CP1$>0$ $\Rightarrow$ alto volumen de cama; CP1$<0$ $\Rightarrow$ servicio especializado/ambulatorio.
\item \textbf{CP2 (Riesgo / hospitalización vs.\ ambulatorio):} cargas positivas en MO, ICM, RE, ES, NE. CP2$>0$ $\Rightarrow$ servicio con alta mortalidad y complejidad.
\end{itemize}

\textbf{(d) Clasificación de los servicios.}
\begin{center}\small
\begin{tabular}{lll}\toprule
Servicio & Cuadrante & Lectura\\\midrule
Medicina Interna & I ($+,+$) & Alta demanda + alta hospitalización\\
Cirugía          & I (cerca origen) & Demanda y hospitalización medias\\
Traumatología, Urología & IV ($+,-$) & Demanda media, ambulatorios\\
Ginecología, Pediatría & IV ($+,-$) & Alta demanda, ambulatorios\\
Digestivo & II/origen & Baja demanda, mortalidad relativa\\
Otorrino, Cardio, Oftalmo & III ($-,-$) & Baja demanda y bajo riesgo\\
Psiquiatría, Hemato, Neuro & II ($-,+$) & Baja demanda, alta complejidad\\
\bottomrule
\end{tabular}\end{center}

\textbf{(e) Mayor carga de trabajo.}
\begin{itemize}[noitemsep,leftmargin=*]
\item \emph{Cuantitativa} (CP1 más positivo): \textbf{Medicina Interna, Ginecología, Pediatría}.
\item \emph{Cualitativa} (CP2 más positivo): \textbf{Medicina Interna} y los del cuadrante II.
\end{itemize}
La que combina ambos es \textbf{Medicina Interna}.

\textbf{(f) Más eficientes.} CP1 alto con CP2 bajo (mucho volumen sin complicaciones) $\Rightarrow$ \textbf{Cirugía y Traumatología}.
\end{solbox}

\begin{figure}[H]
\centering
\begin{tikzpicture}
\begin{axis}[
  width=11cm,height=8cm, axis lines=middle, xlabel={$v_{j1}$ (CP1)}, ylabel={$v_{j2}$ (CP2)},
  xmin=-1.1,xmax=1.1,ymin=-1.1,ymax=1.1, grid=both, grid style={gray!20},
  title={Variables originales P3 (hospital)}]
\draw[gray, dashed] (axis cs:0,0) circle [radius=1];
\arr{0.860}{-0.066}{NI}{south}
\arr{0.421}{0.747}{MO}{west}
\arr{-0.406}{0.670}{RE}{east}
\arr{-0.250}{0.388}{NE}{east}
\arr{-0.562}{0.635}{ICM}{east}
\arr{0.820}{0.508}{ES}{west}
\arr{0.663}{0.078}{IF}{south}
\end{axis}
\end{tikzpicture}
\caption{Variables originales del problema P3.}
\label{fig:p3var}
\end{figure}

\begin{figure}[H]
\centering
\begin{tikzpicture}
\begin{axis}[
  width=12cm,height=9cm, axis lines=middle, xlabel={CP1}, ylabel={CP2},
  xmin=-3,xmax=3.5,ymin=-2.5,ymax=2.5, grid=both, grid style={gray!20},
  title={Servicios hospitalarios (P3)}]
\addplot[only marks,mark=*,mark size=2pt,color=blue]
  coordinates {(2.8,1.7) (-2.5,1.6) (-1.4,1.2) (-1.0,0.7) (0.1,0.2) (1.0,0.5) (1.2,-0.3) (1.5,-0.5) (-0.8,-0.5) (-1.0,-0.7) (-1.0,-1.5) (2.5,-1.6) (1.7,-1.5)};
\node[anchor=south west] at (axis cs:2.8,1.7) {\small Med Int};
\node[anchor=south west] at (axis cs:-2.5,1.6) {\small Psiquiat};
\node[anchor=south west] at (axis cs:-1.4,1.2) {\small Hemato};
\node[anchor=south west] at (axis cs:-1.0,0.7) {\small Neuro};
\node[anchor=south west] at (axis cs:0.1,0.2) {\small Digest};
\node[anchor=south west] at (axis cs:1.0,0.5) {\small Cirugía};
\node[anchor=south west] at (axis cs:1.2,-0.3) {\small Trauma};
\node[anchor=south west] at (axis cs:1.5,-0.5) {\small Urología};
\node[anchor=south west] at (axis cs:-0.8,-0.5) {\small Cardio};
\node[anchor=south west] at (axis cs:-1.0,-0.7) {\small Otorrino};
\node[anchor=south west] at (axis cs:-1.0,-1.5) {\small Oftalmo};
\node[anchor=south west] at (axis cs:2.5,-1.6) {\small Gineco};
\node[anchor=south west] at (axis cs:1.7,-1.5) {\small Pediat};
\end{axis}
\end{tikzpicture}
\caption{Reconstrucción del plano factorial CP1--CP2 con los 13 servicios.}
\end{figure}

\newpage
% =============================================================================
\section{P4. Análisis del sector lechero}
% =============================================================================

\begin{enuncbox}
Haciendas lecheras venezolanas. 6 variables: \textbf{SUP} (superficie), \textbf{VACA} (n.\ vacas), \textbf{SANI} (índice sanitario), \textbf{INST} (instalaciones), \textbf{MAQ} (maquinaria), \textbf{PROM} (leche/vaca).

\smallskip
$\lambda_1=1{,}794$; $\lambda_2=1{,}341$.

\smallskip
\textbf{Cargas factoriales} (autovectores ponderados por $\sqrt{\lambda_i}$):

\begin{center}\small
\begin{tabular}{lrr}\toprule
Variable & CP1 & CP2\\\midrule
SUP  & $\phantom{-}0{,}79$ & $\phantom{-}0{,}10$\\
VACA & $\phantom{-}0{,}76$ & $\phantom{-}0{,}40$\\
SANI & $\phantom{-}0{,}44$ & $-0{,}48$\\
INST & $\phantom{-}0{,}32$ & $-0{,}48$\\
MAQ  & $\phantom{-}0{,}53$ & $-0{,}26$\\
PROM & $-0{,}01$ & $-0{,}80$\\\bottomrule
\end{tabular}\end{center}
\end{enuncbox}

\begin{solbox}
\textbf{Caracterización de componentes.}
\begin{itemize}[noitemsep,leftmargin=*]
\item \textbf{CP1 (Tamaño / dotación productiva):} cargas altas y positivas en SUP, VACA, MAQ. PROM no participa ($-0{,}01$). $\Rightarrow$ ordena las haciendas por \emph{magnitud productiva}.
\item \textbf{CP2 (Tecnificación / industrialización):} cargas \emph{negativas} en PROM, SANI, INST, MAQ. Como los cuatro indicadores de productividad tienen el mismo signo, CP2$<0$ identifica haciendas \emph{tecnificadas}; CP2$>0$ identifica haciendas \emph{extensivas}.
\end{itemize}

\textbf{Identificación de los 8 grupos (A–H).} El enunciado oficial indica que las haciendas se agrupan en 8 tipos:

\begin{center}\small
\begin{tabular}{cl}\toprule
Grupo & Tipo de hacienda\\\midrule
A, C & Mataderos (mucho ganado, baja leche/vaca)\\
B    & Hacienda lechera (alta producción/vaca y tecnificación)\\
D, E & Multipropósito (intermedias)\\
F    & Pastoreo, crianza y engorde (mucha superficie, baja tecnificación)\\
G, H & Haciendas pequeñas (poca superficie y poco ganado)\\\bottomrule
\end{tabular}\end{center}

\textbf{Validez del estudio.} Asumiendo datos estandarizados ($\sum\lambda_i=p=6$):
\[
\mathrm{VE}_2 \;=\; \frac{1{,}794+1{,}341}{6}\cdot 100\% \;=\; \frac{3{,}135}{6}\cdot 100\% \;\approx\; \boxed{52{,}25\%}.
\]
Valor moderado; un análisis riguroso debería incluir CP3 (no entregada en la guía).
\end{solbox}

\begin{figure}[H]
\centering
\begin{tikzpicture}
\begin{axis}[
  width=11cm,height=8cm, axis lines=middle, xlabel={CP1}, ylabel={CP2},
  xmin=-1.1,xmax=1.1,ymin=-1.1,ymax=1.1, grid=both, grid style={gray!20},
  title={Cargas factoriales P4 (sector lechero)}]
\draw[gray, dashed] (axis cs:0,0) circle [radius=1];
\arr{0.79}{0.10}{SUP}{south}
\arr{0.76}{0.40}{VACA}{south west}
\arr{0.44}{-0.48}{SANI}{north}
\arr{0.32}{-0.48}{INST}{north}
\arr{0.53}{-0.26}{MAQ}{east}
\arr{-0.01}{-0.80}{PROM}{east}
\end{axis}
\end{tikzpicture}
\caption{Variables originales del problema P4 (cargas ponderadas).}
\end{figure}

\begin{figure}[H]
\centering
\begin{tikzpicture}
\begin{axis}[
  width=11cm,height=8cm, axis lines=middle, xlabel={CP2}, ylabel={CP1},
  xmin=-3,xmax=3,ymin=-3,ymax=3, grid=both, grid style={gray!20},
  title={8 tipos de haciendas en el plano CP2 (abscisa) -- CP1 (ordenada)}]
\def\pt#1#2#3{\addplot[only marks,mark=diamond*,mark size=3pt,color=blue] coordinates {(#1,#2)}; \node[anchor=south west,font=\small\bfseries] at (axis cs:#1,#2) {#3};}
\pt{-1.8}{1.5}{A}
\pt{-2.5}{0.5}{B}
\pt{-0.8}{2.5}{C}
\pt{-0.3}{0.4}{D}
\pt{0.3}{-0.3}{E}
\pt{1.6}{1.8}{F}
\pt{2.2}{0.05}{G}
\pt{1.4}{-1.5}{H}
\end{axis}
\end{tikzpicture}
\caption{Reconstrucción aproximada de los 8 tipos de haciendas (figura del enunciado).}
\end{figure}

\newpage
% =============================================================================
\section{P5. Análisis de neoplasias}
% =============================================================================

\begin{enuncbox}
Base de datos de tumores (neoplasias) caracterizados por: \textbf{Largo}, \textbf{Ancho}, núcleos \textbf{Regular} e \textbf{Irregulares}. ACP retiene dos componentes.

\begin{center}\small
\begin{tabular}{lrr}\toprule
 & CP1 & CP2\\\midrule
$\lambda$       & $1{,}7034$ & $1{,}2560$\\\midrule
Largo           & $\phantom{-}0{,}5211$ & $-0{,}3774$\\
Regular         & $-0{,}2693$ & $-0{,}9233$\\
Ancho           & $\phantom{-}0{,}5804$ & $-0{,}0245$\\
Irregulares     & $\phantom{-}0{,}5648$ & $-0{,}0669$\\\bottomrule
\end{tabular}\end{center}
\end{enuncbox}

\begin{solbox}
\textbf{(a) Información perdida.} $p=4$ variables estandarizadas, $\sum\lambda_i=4$:
\[
\mathrm{VE}_2 \;=\; \frac{1{,}7034+1{,}2560}{4}\cdot 100\% \;=\; \frac{2{,}9594}{4}\cdot 100\% \;\approx\; 73{,}99\%
\]
\begin{respbox}
$\boxed{\mathrm{IP}_2 \;\approx\; 26{,}01\%}$.
\end{respbox}

\textbf{(b)} Variables en $(v_{j1}, v_{j2})$ — figura~\ref{fig:p5var}.

\textbf{(c) Interpretación.}
\begin{itemize}[noitemsep,leftmargin=*]
\item \textbf{CP1}: cargas positivas en Largo, Ancho e Irregulares y negativa en Regular. Las tres variables que indican mayor tamaño y desorden cargan juntas frente a Regular. $\Rightarrow$ \textbf{CP1 $\equiv$ \emph{Malignidad}} (a mayor CP1, tumor más grande y con núcleos irregulares).
\item \textbf{CP2}: dominada por Regular ($-0{,}92$), con Largo en $-0{,}38$. CP2 muy negativo $\Rightarrow$ alta proporción de núcleos regulares; CP2$\to 0$ o positivo $\Rightarrow$ baja regularidad. $\Rightarrow$ \textbf{CP2 $\equiv$ \emph{Regularidad celular}} (inversa: $-$ regulares, $+$ irregulares).
\end{itemize}

\textbf{(d) ¿Cuántos grupos?} La nube de puntos del enunciado muestra una separación clara en \textbf{dos grupos} a lo largo de CP1: izquierda (CP1$<0$) y derecha (CP1$>0$). Eventualmente puede subdividirse según CP2 en sub-conglomerados, pero \textbf{los grupos principales son 2}.

\textbf{(e) Tumores benignos.} Si malignidad $=$ tamaño $+$ irregularidad $=$ CP1 alto, los \textbf{benignos} corresponden a los tumores con \textbf{CP1$<0$} (lado izquierdo del scatter). En la figura del enunciado, ese cúmulo incluye aproximadamente los puntos numerados \textbf{1--50} (los identificadores menores), por ejemplo: 1, 2, 3, 4, 5, 7--14, 17, 19--50, etc. (cualquier punto con coordenada CP1 negativa). Los puntos con CP1 muy positivo (e.g.\ 119, 120, 131, 136) son los más malignos.
\end{solbox}

\begin{figure}[H]
\centering
\begin{tikzpicture}
\begin{axis}[
  width=11cm,height=8cm, axis lines=middle, xlabel={$v_{j1}$ (CP1)}, ylabel={$v_{j2}$ (CP2)},
  xmin=-1.1,xmax=1.1,ymin=-1.1,ymax=1.1, grid=both, grid style={gray!20},
  title={Variables originales P5 (neoplasias)}]
\draw[gray, dashed] (axis cs:0,0) circle [radius=1];
\arr{0.5211}{-0.3774}{Largo}{north west}
\arr{-0.2693}{-0.9233}{Regular}{east}
\arr{0.5804}{-0.0245}{Ancho}{south}
\arr{0.5648}{-0.0669}{Irregulares}{north west}
\end{axis}
\end{tikzpicture}
\caption{Variables originales del problema P5.}
\label{fig:p5var}
\end{figure}

\begin{figure}[H]
\centering
\begin{tikzpicture}
\begin{axis}[
  width=12cm,height=8cm, axis lines=middle, xlabel={CP1}, ylabel={CP2},
  xmin=-3,xmax=3,ymin=-3,ymax=3, grid=both, grid style={gray!20},
  title={Esquema de la nube de tumores (P5): 2 grupos}]
% Benignos (Q2-Q3, CP1 < 0)
\addplot[only marks,mark=o,mark size=2pt,color=blue,opacity=0.8] coordinates {
  (-2.4,0.2) (-2.0,0.5) (-2.2,-0.4) (-1.8,1.1) (-1.5,-0.7) (-1.3,0.0)
  (-1.0,-1.4) (-1.2,0.9) (-0.9,-1.0) (-0.7,0.4) (-1.6,1.5) (-2.0,-1.0)
  (-0.8,-0.3) (-1.4,-1.2) (-1.1,1.6) (-1.7,-0.2)};
\node[blue,anchor=south] at (axis cs:-1.5,2.0) {\textbf{Benignos}};
% Malignos (CP1 > 0)
\addplot[only marks,mark=triangle*,mark size=2.5pt,color=red!80!black,opacity=0.8] coordinates {
  (0.4,0.9) (0.7,-0.1) (1.0,-0.7) (1.3,0.3) (1.6,-1.0) (1.8,-1.6)
  (2.1,-0.3) (2.4,-1.2) (0.9,-1.4) (1.5,-2.0) (1.2,1.0) (0.6,-2.2)
  (2.0,0.5) (1.7,-0.5) (2.3,-2.0)};
\node[red!80!black,anchor=south] at (axis cs:1.7,1.5) {\textbf{Malignos}};
\draw[dashed,gray] (axis cs:0,-3) -- (axis cs:0,3);
\end{axis}
\end{tikzpicture}
\caption{Esquema cualitativo: separación de tumores benignos vs.\ malignos a lo largo de CP1.}
\end{figure}

\newpage
% =============================================================================
\section{P6. Barómetro Merco}
% =============================================================================

\begin{enuncbox}
50 empresas con más prestigio. Variables: \textbf{REF} (resultados económicos), \textbf{CPS} (calidad producto/servicio), \textbf{CCCL} (cultura corporativa y calidad laboral), \textbf{ERSC} (ética y responsabilidad social), \textbf{DGPI} (dimensión global e internacional), \textbf{IDI} (I+D+i).

\begin{center}\small
\begin{tabular}{lrrrr}\toprule
 & CP1 & CP2 & CP3 & CP4\\\midrule
$\lambda$ & $2{,}15$ & $1{,}32$ & $1{,}18$ & $0{,}77$\\\midrule
REF  & $-0{,}015$ & $\phantom{-}0{,}743$ & $\phantom{-}0{,}296$ & $\phantom{-}0{,}250$\\
CPS  & $\phantom{-}0{,}053$ & $\phantom{-}0{,}066$ & $\phantom{-}0{,}855$ & $-0{,}215$\\
CCCL & $\phantom{-}0{,}582$ & $-0{,}187$ & $\phantom{-}0{,}038$ & $\phantom{-}0{,}439$\\
ERSC & $\phantom{-}0{,}626$ & $-0{,}101$ & $\phantom{-}0{,}207$ & $\phantom{-}0{,}216$\\
DGPI & $\phantom{-}0{,}626$ & $-0{,}101$ & $\phantom{-}0{,}107$ & $\phantom{-}0{,}216$\\
IDI  & $\phantom{-}0{,}245$ & $\phantom{-}0{,}629$ & $-0{,}405$ & $-0{,}036$\\\bottomrule
\end{tabular}\end{center}
\end{enuncbox}

\begin{solbox}
\textbf{(a) Información explicada.} $p=6$, $\sum\lambda_i = 6$.
\begin{align*}
\mathrm{VE}_2 &= \frac{2{,}15+1{,}32}{6}\cdot 100\% = \frac{3{,}47}{6}\cdot 100\% \approx \boxed{57{,}83\%},\\
\mathrm{VE}_3 &= \frac{3{,}47+1{,}18}{6}\cdot 100\% = \frac{4{,}65}{6}\cdot 100\% \approx \boxed{77{,}50\%},\\
\mathrm{VE}_4 &= \frac{4{,}65+0{,}77}{6}\cdot 100\% = \frac{5{,}42}{6}\cdot 100\% \approx \boxed{90{,}33\%}.
\end{align*}

\textbf{(b) Caracterización de cada componente.}
\begin{itemize}[noitemsep,leftmargin=*]
\item \textbf{CP1 — Reputación corporativa y sostenibilidad:} cargas altas en \textbf{CCCL ($0{,}58$)}, \textbf{ERSC ($0{,}63$)}, \textbf{DGPI ($0{,}63$)}; nulas en REF y CPS. Mide \emph{prestigio social, cultural e internacional}.
\item \textbf{CP2 — Desempeño financiero e innovación:} cargas altas en \textbf{REF ($0{,}74$)} e \textbf{IDI ($0{,}63$)}. Mide \emph{resultados económicos y capacidad de innovar}.
\item \textbf{CP3 — Calidad del producto / servicio:} carga muy alta en \textbf{CPS ($0{,}86$)} y negativa en \textbf{IDI ($-0{,}41$)}. Es prácticamente el eje \emph{calidad del producto}, contrapuesto débilmente a innovación.
\item \textbf{CP4 — Cultura laboral residual:} carga moderada en \textbf{CCCL ($0{,}44$)} y \textbf{REF ($0{,}25$)}; el resto bajo. Captura un matiz adicional de cultura organizacional ligado a resultados financieros.
\end{itemize}

\textbf{(c) Identificación de los 8 grupos.} Leyendo simultáneamente los dos planos (CP1--CP2 y CP3--CP4) del enunciado:

\begin{center}\small
\begin{tabular}{cllll}\toprule
Grupo & CP1 & CP2 & CP3 & CP4 \\
      & (reputación) & (econ./I+D) & (calidad) & (cultura)\\\midrule
A & $-{}-$ & $\sim 0$ & $\sim 0$ & $-{}-$ \\
B & $\sim 0$ & $+{}+$  & $\sim 0$ & $-$\\
C & $-$ & $+{}+$  & $+$ & $\sim 0$\\
D & $\sim 0$ & $+$ & $+{}+$ & $\sim 0$\\
E & $\sim 0$ & $\sim 0$ & $\sim 0$ & $-{}-$\\
F & $+$  & $+{}+$  & $-$ & $\sim 0$\\
G & $+{}+$ & $\sim 0$ & $+$ & $\sim 0$\\
H & $\sim 0$ & $-$  & $\sim 0$ & $\sim 0$\\\bottomrule
\end{tabular}\end{center}

\textbf{Interpretación de cada grupo:}
\begin{itemize}[noitemsep,leftmargin=*]
\item \textbf{A}: baja reputación social \emph{y} cultura laboral pobre $\Rightarrow$ empresas peor posicionadas en el ranking.
\item \textbf{B}: muy buenos resultados financieros e innovación pero reputación discreta $\Rightarrow$ empresas \emph{financieramente sólidas}.
\item \textbf{C}: combina buenos resultados/innovación con cierta calidad de producto, con baja reputación social $\Rightarrow$ \emph{tecnológicas emergentes}.
\item \textbf{D}: destaca por \textbf{calidad de producto/servicio} $\Rightarrow$ empresa \emph{centrada en la calidad}.
\item \textbf{E}: posición desfavorable en cultura corporativa $\Rightarrow$ empresa con \emph{problemas laborales internos}.
\item \textbf{F}: reputación alta + resultados económicos altos $\Rightarrow$ \emph{empresa líder global equilibrada}.
\item \textbf{G}: la mejor reputación social/internacional $\Rightarrow$ \emph{marca de prestigio}.
\item \textbf{H}: ligeramente por debajo en innovación $\Rightarrow$ empresa convencional sin atributos sobresalientes.
\end{itemize}
\end{solbox}

\begin{figure}[H]
\centering
\begin{tikzpicture}
\begin{axis}[
  width=11cm,height=8cm, axis lines=middle, xlabel={$v_{j1}$ (CP1)}, ylabel={$v_{j2}$ (CP2)},
  xmin=-1.1,xmax=1.1,ymin=-1.1,ymax=1.1, grid=both, grid style={gray!20},
  title={Variables P6 — plano CP1--CP2}]
\draw[gray, dashed] (axis cs:0,0) circle [radius=1];
\arr{-0.015}{0.743}{REF}{west}
\arr{0.053}{0.066}{CPS}{south west}
\arr{0.582}{-0.187}{CCCL}{east}
\arr{0.626}{-0.101}{ERSC}{south east}
\arr{0.626}{-0.101}{DGPI}{north}
\arr{0.245}{0.629}{IDI}{south west}
\end{axis}
\end{tikzpicture}
\caption{Variables P6 sobre el plano de las dos primeras componentes (Reputación vs.\ Resultados/Innovación).}
\end{figure}

\begin{figure}[H]
\centering
\begin{tikzpicture}
\begin{axis}[
  width=10cm,height=8cm, axis lines=middle, xlabel={CP1}, ylabel={CP2},
  xmin=-1.1,xmax=1.1,ymin=-1.1,ymax=1.1, grid=both, grid style={gray!20},
  title={Grupos de empresas (CP1--CP2)}]
\def\ptt#1#2#3{\addplot[only marks,mark=diamond*,mark size=3pt,color=usachteal!80!black] coordinates {(#1,#2)}; \node[anchor=south west,font=\small\bfseries] at (axis cs:#1,#2) {#3};}
\ptt{-0.65}{0.0}{A}
\ptt{0.05}{0.85}{B}
\ptt{-0.10}{0.75}{C}
\ptt{-0.05}{0.18}{D}
\ptt{0.0}{-0.05}{E}
\ptt{0.40}{0.60}{F}
\ptt{0.85}{0.05}{G}
\ptt{0.10}{-0.20}{H}
\end{axis}
\end{tikzpicture}
\hfill
\begin{tikzpicture}
\begin{axis}[
  width=10cm,height=8cm, axis lines=middle, xlabel={CP3}, ylabel={CP4},
  xmin=-1.1,xmax=1.1,ymin=-1.1,ymax=1.1, grid=both, grid style={gray!20},
  title={Grupos de empresas (CP3--CP4)}]
\def\pto#1#2#3{\addplot[only marks,mark=diamond*,mark size=3pt,color=orange!80!black] coordinates {(#1,#2)}; \node[anchor=south west,font=\small\bfseries] at (axis cs:#1,#2) {#3};}
\pto{0.05}{-0.75}{A}
\pto{0.05}{-0.05}{B}
\pto{0.15}{-0.05}{C}
\pto{0.80}{0.05}{D}
\pto{0.05}{-0.85}{E}
\pto{-0.20}{0.05}{F}
\pto{0.10}{0.10}{G}
\pto{0.0}{0.0}{H}
\end{axis}
\end{tikzpicture}
\caption{Proyección esquemática de los 8 grupos de empresas en ambos planos factoriales del Barómetro Merco.}
\end{figure}

\newpage
% =============================================================================
\section{P7. Clasificación de neuronas}
% =============================================================================

\begin{enuncbox}
6 tipos de neuronas (A--F) caracterizadas por: \textbf{Cantidad D} (n.\ de dendritas), \textbf{Arborización D}, \textbf{Largo Axón}, \textbf{Tamaño Soma}. ACP a 3 componentes:

\begin{center}\small
\begin{tabular}{lrrrrr}\toprule
 & $\lambda$ & Cantidad D & Arborización D & Largo Axón & Tamaño Soma\\\midrule
CP1 & $1{,}8$ & $\phantom{-}0{,}948$ & $\phantom{-}0{,}925$ & $\phantom{-}0{,}769$ & $\phantom{-}0{,}142$\\
CP2 & $0{,}9$ & $-0{,}215$ & $\phantom{-}0{,}311$ & $\phantom{-}0{,}629$ & $-0{,}171$\\
CP3 & $0{,}3$ & $\phantom{-}0{,}005$ & $-0{,}004$ & $-0{,}007$ & $\phantom{-}0{,}789$\\\bottomrule
\end{tabular}\end{center}

Coordenadas de los 6 tipos en CP1--CP2:

\begin{center}\small
\begin{tabular}{lrrrrrr}\toprule
 & A & B & C & D & E & F\\\midrule
CP1 & $-0{,}60$ & $-1{,}10$ & $-0{,}51$ & $0{,}40$ & $1{,}20$ & $1{,}30$\\
CP2 & $\phantom{-}1{,}60$ & $\phantom{-}0{,}01$ & $\phantom{-}0{,}81$ & $0{,}95$ & $-0{,}55$ & $0{,}63$\\\bottomrule
\end{tabular}\end{center}
\end{enuncbox}

\begin{solbox}
\textbf{Validez a dos componentes.} $p=4$. Como se entregan tres autovalores ($1{,}8+0{,}9+0{,}3=3{,}0$), el cuarto debe valer $\lambda_4 = 4-3 = 1{,}0$:
\[
\mathrm{VE}_2 \;=\; \frac{1{,}8+0{,}9}{4}\cdot 100\% \;=\; \frac{2{,}7}{4}\cdot 100\% \;=\; \boxed{67{,}5\%}.
\]
Información perdida: $\boxed{32{,}5\%}$.

\textbf{Variables en el plano.} Cada variable en $(v_{j1},v_{j2})$ — figura~\ref{fig:p7var}.

\textbf{Interpretación de componentes.}
\begin{itemize}[noitemsep,leftmargin=*]
\item \textbf{CP1 — Complejidad dendrítica--axonal:} las tres variables Cantidad D ($0{,}95$), Arborización D ($0{,}93$) y Largo Axón ($0{,}77$) cargan altísimo y positivo. Tamaño Soma es casi nulo. $\Rightarrow$ CP1$>0$ indica neurona \emph{más desarrollada} (más dendritas, más ramificada, axón largo).
\item \textbf{CP2 — Axón largo vs.\ muchas dendritas:} Largo Axón ($+0{,}63$) y Arborización D ($+0{,}31$) frente a Cantidad D ($-0{,}22$). $\Rightarrow$ CP2 separa neuronas \emph{tipo proyección} (axón largo, pocas dendritas, CP2$>0$) de neuronas \emph{tipo interneurona} (muchas dendritas, CP2$<0$).
\item \textbf{CP3 — Tamaño del soma:} prácticamente sólo Tamaño Soma carga ($0{,}79$). Es un eje específico para el cuerpo celular.
\end{itemize}

\textbf{Caracterización de los 6 tipos.}
\begin{center}\small
\begin{tabular}{cccp{8cm}}\toprule
Tipo & CP1 & CP2 & Lectura\\\midrule
A & $-0{,}60$ & $\phantom{-}1{,}60$ & Poco desarrollada pero con axón muy largo $\Rightarrow$ neurona \emph{de proyección}, simple, con axón dominante.\\
B & $-1{,}10$ & $\phantom{-}0{,}01$ & La menos desarrollada: pocas dendritas, sin axón largo. Neurona \emph{pequeña / inmadura}.\\
C & $-0{,}51$ & $\phantom{-}0{,}81$ & Poco desarrollada con tendencia a axón largo. Similar a A pero más moderada.\\
D & $\phantom{-}0{,}40$ & $\phantom{-}0{,}95$ & Desarrollo medio + axón largo. \emph{Piramidal de proyección}.\\
E & $\phantom{-}1{,}20$ & $-0{,}55$ & Muy desarrollada (muchas dendritas y ramificación) pero sin axón largo $\Rightarrow$ \emph{interneurona local} muy ramificada.\\
F & $\phantom{-}1{,}30$ & $\phantom{-}0{,}63$ & La más desarrollada en todo $\Rightarrow$ neurona \emph{piramidal grande} (mucha arborización + axón largo).\\\bottomrule
\end{tabular}\end{center}
\end{solbox}

\begin{figure}[H]
\centering
\begin{tikzpicture}
\begin{axis}[
  width=11cm,height=8cm, axis lines=middle, xlabel={$v_{j1}$ (CP1)}, ylabel={$v_{j2}$ (CP2)},
  xmin=-1.1,xmax=1.1,ymin=-1.1,ymax=1.1, grid=both, grid style={gray!20},
  title={Variables originales P7 (neuronas)}]
\draw[gray, dashed] (axis cs:0,0) circle [radius=1];
\arr{0.948}{-0.215}{Cant.\ D}{north}
\arr{0.925}{0.311}{Arboriz.\ D}{south}
\arr{0.769}{0.629}{Largo Axón}{south west}
\arr{0.142}{-0.171}{Tamaño Soma}{north}
\end{axis}
\end{tikzpicture}
\caption{Variables del problema P7 proyectadas en el plano CP1--CP2.}
\label{fig:p7var}
\end{figure}

\begin{figure}[H]
\centering
\begin{tikzpicture}
\begin{axis}[
  width=12cm,height=8cm, axis lines=middle, xlabel={CP1}, ylabel={CP2},
  xmin=-2,xmax=2,ymin=-2,ymax=2.2, grid=both, grid style={gray!20},
  title={Los 6 tipos de neuronas en el plano CP1--CP2}]
\def\ptn#1#2#3{\addplot[only marks,mark=*,mark size=3pt,color=red!70!black] coordinates {(#1,#2)}; \node[anchor=south west,font=\small\bfseries] at (axis cs:#1,#2) {#3};}
\ptn{-0.60}{1.60}{A}
\ptn{-1.10}{0.01}{B}
\ptn{-0.51}{0.81}{C}
\ptn{0.40}{0.95}{D}
\ptn{1.20}{-0.55}{E}
\ptn{1.30}{0.63}{F}
\end{axis}
\end{tikzpicture}
\caption{Distribución de los 6 tipos de neuronas en el plano factorial principal.}
\end{figure}

\newpage
% =============================================================================
\section{Resumen comparativo de los 7 problemas}
% =============================================================================

\begin{center}\small
\begin{tabular}{clcccp{4cm}p{4cm}}\toprule
\# & Problema & $p$ & $\mathrm{VE}_2$ & $\mathrm{IP}_2$ & Interpretación CP1 & Interpretación CP2\\\midrule
P1 & Consumo familias    & 7 & $78{,}4\%$ & $21{,}6\%$ & Calidad de canasta (proxy ingreso) & Estilo de dieta (lácteos+pan vs.\ proteína)\\
P2 & Billetes            & 6 & $65{,}3\%$ & $34{,}7\%$ & Forma / márgenes              & Tamaño longitudinal\\
P2 & Monedas             & 6 & $58{,}3\%$ & $41{,}7\%$ & ---                            & ---\\
P3 & Hospital            & 7 & $62{,}7\%$ & $37{,}3\%$ & Demanda / volumen              & Hospitalización vs.\ ambulatorio\\
P4 & Sector lechero      & 6 & $52{,}3\%$ & $47{,}7\%$ & Tamaño / dotación             & Tecnificación (PROM/SANI)\\
P5 & Neoplasias          & 4 & $74{,}0\%$ & $26{,}0\%$ & Malignidad (tamaño+irregul.)   & Regularidad celular\\
P6 & Barómetro Merco     & 6 & $57{,}8\%$ & $42{,}2\%$ & Reputación social/global       & Resultados económicos + I+D\\
P7 & Neuronas            & 4 & $67{,}5\%$ & $32{,}5\%$ & Complejidad dendrítica-axonal  & Axón largo vs.\ muchas dendritas\\
\bottomrule
\end{tabular}\end{center}

\bigskip
\textbf{Observaciones metodológicas finales.}
\begin{itemize}[leftmargin=*]
\item Todos los problemas asumen \emph{datos estandarizados} ($\sum\lambda_i = p$). En P4 y P7 se completaron autovalores faltantes bajo este supuesto.
\item En P2 se corrigió la respuesta de monedas a $58{,}33\%$ (la pauta original reporta $58{,}83\%$, probable errata).
\item En P5 la asignación numérica precisa de tumores benignos depende de la figura exacta del enunciado; la regla operativa es \textbf{tumor benigno $\Leftrightarrow$ CP1$<0$}.
\item En P6 el rótulo de los grupos se infiere de los dos biplots (CP1--CP2 y CP3--CP4) entregados; una asignación exacta requeriría las coordenadas numéricas.
\end{itemize}

\end{document}
